\documentclass[10pt,a4paper]{amsart}
\usepackage[margin=19mm]{geometry}
\usepackage{amsmath,amssymb,mathtools}
\usepackage[scr=rsfs,cal=euler]{mathalfa}
\usepackage{xfrac}
\usepackage[unicode,hidelinks,bookmarks=false]{hyperref}
\newcommand{\ie}{i.e.\ }
\newcommand{\cf}{cf.\ }
\newcommand{\resp}{respectively\ }
\newcommand{\Subsection}{\subsection}
\providecommand{\nobreakdash}{\nobreak\hbox{-}}
\setlength{\emergencystretch}{2em}
\multlinegap0pt
\usepackage{bm}
\usepackage{bbm}
\usepackage{dsfont}
\usepackage{xspace}
\usepackage{cancel}
\usepackage{mathtools}
\usepackage{amsthm}
\makeatletter
\@ifundefined{lemma}{\newtheorem{lemma}{Lemma}}{}
\makeatother
\newcommand{\Ycal}{\mathcal{Y}}
\title[Is Zero-Shot Super-Resolution Possible in Operator Learning?]{Is Zero-Shot Super-Resolution Possible in Operator Learning?}
\author{Unique Subedi, Ambuj Tewari}
\date{Source: arXiv:2606.00296v1 (CC BY 4.0)}
\begin{document}
\maketitle

\noindent\emph{Source and licence.} Is Zero-Shot Super-Resolution Possible in Operator Learning?. Version arXiv:2606.00296v1 (CC BY 4.0), distributed under the Creative Commons Attribution 4.0 International licence, as recorded in the arXiv licence metadata for this exact version and shown on the abstract page https://arxiv.org/abs/2606.00296v1. Full rights notice: licenses/arxiv-2606.00296-CC-BY-4.0.txt.\par\smallskip

\noindent\textit{Editorial scope.} Counted unit: the Lemma whose source label is \texttt{None} in the article (source0.tex lines 1264--1271, source label \texttt{None}), delivered together with the article's own complete original proof of it (source0.tex lines 1273--1337, one \texttt{proof} environment of 65 lines). No mathematics is written, repaired or re-derived: statement and proof are the article's own text line for line. Required context retained uncounted: none. Every definition, display and label of those lines is retained unchanged. Omitted from this package and not redistributed: 1--1263, 1272--1272, 1338--1435. Nothing inside a retained excerpt is omitted or altered: every retained body is delivered exactly as the article's lines are written, so no omission receipt of any kind accompanies this package. Citations: the retained excerpts carry no citation command at all, so no bibliography entry of the article is transcribed and no reference is redistributed. Reference adaptation: none was needed and none was made; every reference inside the delivered unit resolves inside it, so no unresolved reference or citation is produced. Printed slips kept literal: no printed word, symbol, hypothesis or number of the counted statement or of its proof was altered. Layout and class: the article's own class is \texttt{article} with packages including inputenc, amsmath, amsfonts, amssymb, amsthm, graphicx, caption, xcolor, float, hyperref, natbib, verbatim, mathtools, tikz; the wrapper is the lane's pinned \texttt{amsart} editorial wrapper at 10pt on A4, which restates the theorem-like environments the retained text needs and adds the article's own macro definitions that the retained text needs (\texttt{\textbackslash Ycal}), with extra wrapper packages (bm, bbm, dsfont, xspace, cancel, mathtools). The delivered numbering is the unit's own. Assets: no retained span contains a figure environment, a graphics-inclusion command, a PDF-page inclusion, a table or a TikZ picture; no image file is copied into this package, the package declares no asset dependency and no image file is redistributed with it. Licence: version arXiv:2606.00296v1 of this article is distributed under the Creative Commons Attribution 4.0 International licence, as recorded in the arXiv licence metadata for this exact version and shown on the abstract page https://arxiv.org/abs/2606.00296v1. A full rights notice is delivered at licenses/arxiv-2606.00296-CC-BY-4.0.txt. Characters: every retained excerpt is pure ASCII, so the delivered engine, XeLaTeX with its default Latin Modern text font, reports no missing character.
\begin{lemma}
Let $\Ycal_k \subseteq \Ycal_\ell \subseteq \Ycal$ be finite sets. Let $h_k$ denote the fill distance of $\Ycal_k$ and $q_\ell$ the separation distance of $\Ycal_\ell$. Then
\[
\nu_{k,\ell}
\le
\left(1+\frac{h_k}{q_\ell}\right)^d.
\]
\end{lemma}

\begin{proof}
Fix $z \in \Ycal_k$ and define
\[
A_z := \{y \in \Ycal_\ell : \mathrm{nn}(y)=z\}.
\]
We will bound $|A_z|$ uniformly in $z$. If $y \in A_z$, then by definition of nearest neighbor,
\[
|y - z|_2 = \min_{z' \in \Ycal_k} |y - z'|_2 \le h_k.
\]
Hence,
\[
A_z \subseteq B(z,h_k).
\]

By definition of the separation distance $q_\ell$, any two distinct points $y,y' \in \Ycal_\ell$ satisfy
\[
|y - y'|_2 \ge 2q_\ell.
\]
Therefore, the open balls $\{B^\circ(y,q_\ell)\}_{y \in A_z}$ are pairwise disjoint.

Let $y \in A_z$ and $x \in B^\circ(y,q_\ell)$. Then
\[
|x - z|_2 \le |x - y|_2 + |y - z|_2 < q_\ell + h_k,
\]
so
\[
B^\circ(y,q_\ell) \subseteq B(z,h_k + q_\ell)
\quad \text{for all } y \in A_z.
\]
Since this is true for all $y$ and $z$, we have
\[
\bigcup_{y \in A_z} B^\circ(y,q_\ell) \subseteq B(z,h_k + q_\ell).
\]
Since the sets $\{B^\circ(y,q_\ell)\}_{y \in A_z}$ are pairwise disjoint,
\[
\operatorname{vol}\!\left(\bigcup_{y \in A_z} B^\circ(y,q_\ell)\right)
=
\sum_{y \in A_z} \operatorname{vol}(B^\circ(y,q_\ell))
=
|A_z| \cdot \operatorname{vol}(B^\circ(0,q_\ell)),
\]
where the second equality holds due to translation invariance of Lebesgue measure.
Since the union is contained in $B(z,h_k + q_\ell)$, we obtain
\[
|A_z| \cdot \operatorname{vol}(B^\circ(0,q_\ell))
\le
\operatorname{vol}(B(z,h_k + q_\ell)).
\]
Using translation invariance of Lebesgue measure and the fact that open and closed balls have the same volume,
\[
|A_z| \cdot \operatorname{vol}(B(0,q_\ell))
\le
\operatorname{vol}(B(0,h_k + q_\ell)).
\]
Finally, using the scaling of volume in $\mathbb{R}^d$, we conclude
\[
|A_z|
\le
\left(\frac{h_k + q_\ell}{q_\ell}\right)^d
=
\left(1+\frac{h_k}{q_\ell}\right)^d.
\]

Taking the maximum over $z \in \Ycal_k$ yields the desired bound on $\nu_{k,\ell}$.
\end{proof}

\end{document}
